\section{Introduction}\label{sec:introduction}

For a finite set \(X\subset\R^2\), write
\[
 u(X)=\#\bigl\{\{x,y\}\subset X:\ \norm{x-y}_2=1\bigr\},
 \qquad
 u(n)=\max_{\substack{X\subset\R^2\\ |X|=n}}u(X).
\]
Thus pairs are unordered and their endpoints are distinct. We prove a fixed
power saving over the exponent \(4/3\).

\begin{theorem}\label{thm:main}
There are absolute constants \(0<C<\infty\) and \(1\le\beta<4/3\) such that
\(u(n)\le Cn^\beta\) for every nonnegative integer \(n\).
\end{theorem}

\Cref{thm:main} gives a positive answer to the question whether this upper
exponent can be lowered by a fixed amount: equivalently,
\(u(n)=O(n^{4/3-\delta})\) for an absolute \(\delta>0\).
All constants are independent of the configuration.

\subsection{The unit-distance problem and earlier methods}

Erd\H{o}s introduced the repeated-distance problem in 1946
\cite{erdos1946}. Rescaling a configuration turns any specified positive
distance into a unit distance, so the problem concerns the largest possible
multiplicity of one distance. His lattice construction gives the lower bound
\(n^{1+c/\log\log n}\) for a positive constant \(c\) and sufficiently large
\(n\) \cite[Theorem~2]{erdos1946}.
Spencer, Szemer\'edi and Trotter established the upper bound
\(u(n)=O(n^{4/3})\) in 1984 \cite{sst1984}.
Sz\'ekely gave a crossing-number proof and a general incidence theorem
for families of curves \cite[Theorems~4 and~8]{Szekely1997}.
The crossing inequality was developed independently by Ajtai,
Chv\'atal, Newborn and Szemer\'edi and by Leighton
\cite{ajtaiEtAl1982,leighton1983}.
Sz\'ekely's incidence theorem applies to points and unit circles using two elementary
properties: two unit circles meet in at most two points, and two points
belong to at most two unit circles.
Later work improved the leading constant
\cite{agostonPalvolgyi2022} and described the structure of configurations
near the extremal incidence scale \cite{katzSilier2026}.
Pach, Raz and Solymosi connected the problem to rigidity and obtained
a logarithmic improvement conditional on a rigidity conjecture
\cite[Conjecture~7 and Theorem~8]{pachRazSolymosi2026}.

The historical unit-distance conjecture of Erd\H{o}s proposed the stronger
bound \(u(n)=n^{1+o(1)}\). An OpenAI construction disproved it by giving
at least \(n^{1+\varepsilon}\) unit distances, for some fixed
\(\varepsilon>0\), along an unbounded sequence of sizes
\cite[Theorem~1.1]{openaiUnitLower2026}. Alon et al.\ gave a simplified
account and explained the construction's number-theoretic ancestry
\cite[Sections~1.2--1.3]{alonEtAl2026}. They credit the class-group
ideas to Michel, Soundararajan, Ellenberg and Venkatesh. The construction
combines these ideas with Golod--Shafarevich
number-field towers having prescribed splitting, including the
tower-cutting method of Hajir, Maire and Ramakrishna
\cite{ellenbergVenkatesh2007,golodShafarevich1964,hajirMaireRamakrishna2021}.
Sawin sharpened the construction to \(u(n)\gg n^{1.014114}\) along an
unbounded sequence of sizes \cite[Theorem~1]{sawin2026}.
These are construction lower bounds. Theorem~\ref{thm:main} bounds every
configuration from above, leaving a gap between the known exponents.

Our proof begins with the same point--circle incidence interpretation.
Random cuttings isolate a piece of the incidence graph with controlled
degrees; their sampled-cell construction belongs to the framework of
Clarkson and Clarkson--Shor \cite{Clarkson1987,ClarksonShor1989}.
The next ingredient is information: two points sampled as neighbors of
one center share little information. Entropy methods already played a
role in the distinct-distance estimates of Tardos and Katz--Tardos
\cite{Tardos2003,KatzTardos}. In particular, the conditional-copy
argument of Katz--Tardos retains the marginal law when resampling
within a class. Here that principle enters a short-history prediction
lemma whose tests may depend on a neighboring center. Its exact
conditional laws and the independence of the fresh test are proved in
\Cref{sec:prediction}.

The arithmetic input is the normalized product formula
\cite{ProductFormula}. We use the associated absolute height of
algebraic numbers: the weighted sum of the positive parts of their
logarithmic absolute values. The product formula converts
relations between unit-edge jumps into bounds summed over all absolute
values. The proof then combines these bounds with prediction and, at its
last stage, algebraic independence and quadratic extensions. The
intermediate statements connecting probability, height and algebra are
proved in the paper.

\subsection{The proof in outline}

Suppose that no fixed power saving exists. Set \(t=n^{1/3}\) along a sequence
of counterexamples, so their ordered unit pairs number \(t^{4-o(1)}\).
The following steps produce a contradiction.

\paragraph{An incidence graph and its sampling laws.}
Use two copies of the configuration as points and unit-circle centers.
An edge joins a point to a center at unit distance. Algebraic
specialization preserves these equations and distinctness. In
\Cref{sec:incidence} we restrict this graph to a piece with controlled
degrees and an expansion estimate that limits the fraction of edges
captured by small vertex subsets. Throughout, a random edge is uniform
on the retained graph, so its vertex marginals are proportional to
degrees. A \emph{wedge} is sampled by choosing a center with this
marginal and then two of its neighbors independently. Its endpoint
mutual information---the relative entropy of their joint law from the
product of their marginals---is \(o(\log t)\).

\paragraph{Prediction from a short history.}
A test assigns a finite code to each point. A history records its codes
at a short list of previous tests. Given an original point, sample a
prediction from the point marginal conditional on having the same
history, and sample a neighboring center from the original point.
The prediction keeps the original point marginal. The information
bound allows a history to be fixed for which the prediction's answers
are sufficiently compatible with the sampled center at a fresh test,
independent of this whole experiment. Section~\ref{sec:prediction}
proves this statement for general finite point--center laws and then
adapts it to the coordinate tests used below.

\paragraph{Arithmetic scale and the bounded case.}
In the complex coordinates \(z=x+iy\) and \(w=x-iy\), a unit edge
has jumps \(d\) and \(d^{-1}\). At an absolute value \(v\) and a real
threshold \(s\), record the bit \(\1_{\{-\log|d|_v>s\}}\).
Let \(f(v,s)\) be the smaller of the fractions of edges with bit zero
and bit one. The scale \(S\) is the integral of \(f\) over all
thresholds, summed with the normalized place weights. Thus \(S\)
measures how much the logarithmic edge-jump profiles vary.

If \(S\) stays bounded, the product formula transfers this control
to heights of normalized point coordinates. At a fresh complex
embedding, quantization tests produce a predicted point separated
from the original point, but approximately satisfying the unit
equations at three of its neighbors. The three normalized jumps are
also separated in that embedding. A determinant calculation excludes
this configuration. Section~\ref{sec:heights} proves the required
height bounds and this contradiction; the remaining case is
\(S\to\infty\).

\paragraph{Common and rare levels.}
The same thresholds give each vertex two labels, recording closeness
in its two coordinates; their pair is called its state.
Except for an integrated error of \(O(1)\),
exactly one label agrees across an edge, and the bit tells us which one.
Call levels common when both bits have appreciable frequency, and
rare when one bit has small frequency. The rare levels with large
consistency error have total length \(o(S)\) and can be removed.
On the remaining rare levels one coordinate label occurs at almost
every vertex; its exceptions are measured by their total level length.

For two edges sharing a point or sharing a center, assign scores that
reward certain matches of the other endpoints and charge bit
disagreements. The product formula bounds their integrals in
\Cref{sec:budgets}. The prediction lemma and the elementary geometry
of two coordinate labels give lower bounds on common levels, while
the rare-level terms cancel when the two orientations are added.
One step requires care: concentration of the predicted states on a
row (a fixed first label), a column (a fixed second label), or at most
two states must be transferred to the actual point, even when that
support is selected using sampled histories. Lemma~\ref{com:transfer}
proves this comparison under the actual edge and history laws.
Section~\ref{sec:common} concludes that common levels contribute
negligibly to center-wedge variation, but a positive amount of
point-exception mass remains on rare levels.

\paragraph{Height-weighted extraction.}
The remaining mass provides many endpoint pairs \((p,r)\) of wedges,
each with a chosen center \(q\) and ratio
\(\lambda_{pr}=(z_p-z_q)/(z_r-z_q)\).
The logarithms of these ratios are approximately differences of
profiles attached to the two endpoints. Consequently, on a complete
rectangle of retained pairs, the alternating product
\[
 \frac{\lambda_{pr}\lambda_{p'r'}}{\lambda_{pr'}\lambda_{p'r}}
\]
has small height. Section~\ref{sec:extraction} produces a dense graph
where all individual ratios have height at least \(J/3\), with
\(J\to\infty\), and every rectangle product has height \(o(J)\).
The selection first removes errors in their height-weighted measure,
then chooses a common height band: small exceptional probability alone
would not control lost height. Pairs on any unit circle containing
many points are removed as well.

\paragraph{An algebraic obstruction.}
Section~\ref{sec:algebra} rules out this dense pair graph. In a field
formed from the sequence, elements of height \(o(J)\) form a
relatively algebraically closed subfield. The rectangle products lie
in this subfield, while the individual ratios are transcendental over
it. We select a fixed complete bipartite pattern on plane or curve
loci. Over the field generated by the defining coefficients of these
loci, its positions have no polynomial relations beyond those required
by the loci. Differentiation over the negligible-height subfield
then gives a finite system involving radicals derived from the
pair-distance identity.

After adjoining one side's positions and derivative data, a point on
the other side still has the independence required by its locus.
For each radicand we construct, over this enlarged
coefficient field, a valuation that is odd on that radicand and
zero on the others. These valuations permit independent sign changes
of the square roots. An unused sign change contradicts the
differentiated equations. The popular-unit-circle exclusion removes
the case in which the required square-class independence fails.

The proof excludes every sequence whose unit-distance counts approach
the exponent \(4/3\), establishing the existence of a fixed positive
gap without estimating its size.

Throughout the proof, asymptotic notation refers to the chosen sequence
\(t\to\infty\). Constants denoted by \(c,C\) may change between occurrences.
All logarithms are natural. Every appeal to a slowly varying parameter
will be made after the quantities it must dominate have been specified.
